\subsubsection{Stride}\label{sec:stride}

In short, \definition{Stride} \hl{controls how far the filter moves between two consecutive convolution positions.} With stride $1$, the filter moves one pixel at a time. With stride $2$, it ``\emph{jumps}'' every second row and every second column. Larger stride values reduce the spatial extent of the output, and that stride $>1$ is conceptually equivalent to convolution followed by downsampling, but without computing values that would then be discarded (\emph{optimization}).

\highspace
Suppose we have a one-dimensional input first. With stride $S=1$, the kernel moves as: position 0, position 1, position 2, position 3, and so on. With stride $S=2$, the kernel moves as: position 0, position 2, position 4, and so on. So we compute \textbf{fewer output activations} with larger stride values. That is why stride \hl{performs a kind of spatial downsampling (i.e., reduces the spatial dimensions of the output).}

\highspace
\textcolor{Green3}{\faIcon{question-circle} \textbf{Why stride is downsampling?}} Because it skips some input locations, and therefore the output is smaller. Imagine first doing the full stride-1 convolution:
\begin{equation*}
    [a_0, a_1, a_2, a_3, a_4]
\end{equation*}
Then keep only every second value:
\begin{equation*}
    [a_0, a_2, a_4]
\end{equation*}
That gives the same sampling pattern as a convolution performed directly with stride $2$. So conceptually:
\begin{equation*}
    \text{stride-2 convolution} \approx \text{stride-1 convolution} + \text{downsampling by 2}
\end{equation*}
Notice that \hl{doing the two operations separately would be \textbf{less efficient} because many intermediate activations would be computed and then discarded.} Instead, stride $>1$ allows us to skip those computations entirely.

\highspace
\begin{examplebox}[: Stride and Downsampling]
    Suppose:
    \begin{equation*}
        N_{\text{in}} = 7, \quad K = 3, \quad P = 0
    \end{equation*}
    With stride $1$, the valid kernel starting position are:
    \begin{equation*}
        0, 1, 2, 3, 4
    \end{equation*}
    So the output size is $N_{\text{out}} = 5$. With stride $2$, instead, we only use:
    \begin{equation*}
        0, 2, 4
    \end{equation*}
    Therefore $N_{\text{out}} = 3$. Same input, same kernel, fewer positions because we skip locations.
\end{examplebox}

\newpage

\begin{figure}[htbp]
\centering
\tikzsetnextfilename{stride-convolution-comparison-1-vs-2}
\tikzwidth{\textwidth}
\begin{tikzpicture}[
    cell/.style={rectangle, draw=black!50, minimum size=1.1cm, inner sep=0pt, anchor=center, font=\small},
    in cell/.style={cell, fill=blue!8, draw=blue!60!black},
    skip cell/.style={cell, fill=black!5, draw=black!30, font=\color{black!40}\small},
    out cell/.style={cell, fill=green!10, draw=green!60!black},
    header/.style={font=\bfseries, align=center},
    labeltext/.style={font=\small, align=center}
]

% =========================================================================
% TOP PANEL: STRIDE S = 1 (UNIT STEP)
% =========================================================================
\begin{scope}[local bounding box=scope1]

    % Panel Title
    \node[header, anchor=south] at (3.6, 5.3) {1. Stride $S=1$ (Move 1 position at a time)};

    % Input Strip 1x7
    \foreach \x in {0,...,6} {
        \node[in cell] (s1-in-\x) at (\x*1.2, 3) {$x_\x$};
    }
    \node[font=\footnotesize\bfseries, color=blue!80!black, anchor=east] at (-0.8, 3) {Input ($N_{\text{in}}=7$):};

    % First Filter Position [0..2] highlighted
    \filldraw[draw=orange!90!black, fill=orange!25, fill opacity=0.45, line width=1.2pt, rounded corners=2pt]
        (s1-in-0.north west) rectangle (s1-in-2.south east);

    % Consecutive step arc arrows showing movement by +1
    \draw[{Stealth[length=3.5pt]}-{Stealth[length=3.5pt]}, orange!90!black, line width=0.9pt]
        (s1-in-1.north) to[out=50, in=130] node[above=-1pt, font=\tiny\bfseries] {+1} (s1-in-2.north);
    \draw[{Stealth[length=3.5pt]}-{Stealth[length=3.5pt]}, orange!90!black, line width=0.9pt]
        (s1-in-2.north) to[out=50, in=130] node[above=-1pt, font=\tiny\bfseries] {+1} (s1-in-3.north);
    \draw[{Stealth[length=3.5pt]}-{Stealth[length=3.5pt]}, orange!90!black, line width=0.9pt]
        (s1-in-3.north) to[out=50, in=130] node[above=-1pt, font=\tiny\bfseries] {+1} (s1-in-4.north);
    \draw[{Stealth[length=3.5pt]}-{Stealth[length=3.5pt]}, orange!90!black, line width=0.9pt]
        (s1-in-4.north) to[out=50, in=130] node[above=-1pt, font=\tiny\bfseries] {+1} (s1-in-5.north);

    % Arrow down to output
    \draw[->, >=stealth, line width=1pt, black!60] (3.6, 2.2) -- (3.6, 1.2)
        node[midway, right, font=\scriptsize] {$K=3, P=0$};

    % Output Strip 1x5
    \foreach \x in {1,...,5} {
        \pgfmathsetmacro{\idx}{int(\x-1)}
        \ifnum\x=1
            \node[out cell, fill=orange!40, draw=orange!80!black, line width=1.2pt] (s1-out-\idx) at (\x*1.2, 0.4) {$a_\idx$};
        \else
            \node[out cell] (s1-out-\idx) at (\x*1.2, 0.4) {$a_\idx$};
        \fi
    }
    \node[font=\footnotesize\bfseries, color=green!50!black, anchor=east] at (-0.8, 0.4) {Output ($N_{\text{out}}=5$):};

    % Bottom Formula
    \node[labeltext, anchor=north] at (3.6, -0.55) {
        $N_{\text{out}} = N_{\text{in}} + 2P - K + 1 = 7 + 0 - 3 + 1 = \mathbf{5}$ valid activations
    };
\end{scope}

% =========================================================================
% SEPARATOR LINE
% =========================================================================
\draw[dashed, black!25, line width=0.8pt] (-1.5, -1.75) -- (8.7, -1.75);

% =========================================================================
% BOTTOM PANEL: STRIDE S = 2
% =========================================================================
\begin{scope}[shift={(0, -7.8)}, local bounding box=scope2]

    % Panel Title
    \node[header, anchor=south] at (3.6, 5.3) {2. Stride $S=2$ (Jump 2 positions at a time / Downsampling)};

    % Input Strip 1x7
    \foreach \x in {0,...,6} {
        \node[in cell] (s2-in-\x) at (\x*1.2, 3) {$x_\x$};
    }
    \node[font=\footnotesize\bfseries, color=blue!80!black, anchor=east] at (-0.8, 3) {Input ($N_{\text{in}}=7$):};

    % First Filter Position [0..2] highlighted
    \filldraw[draw=orange!90!black, fill=orange!25, fill opacity=0.45, line width=1.2pt, rounded corners=2pt]
        (s2-in-0.north west) rectangle (s2-in-2.south east);

    % Jump arc arrows showing movement by +2
    \draw[-{Stealth[length=4pt]}, orange!90!black, line width=1.1pt]
        (s2-in-1.north) to[out=55, in=125] node[above=0pt, font=\tiny\bfseries] {+2 (skip 1)} (s2-in-3.north);
    \draw[-{Stealth[length=4pt]}, orange!90!black, line width=1.1pt]
        (s2-in-3.north) to[out=55, in=125] node[above=0pt, font=\tiny\bfseries] {+2 (skip 3)} (s2-in-5.north);

    % Highlight valid starting centers (0, 2, 4)
    \node[circle, draw=orange!90!black, fill=orange!40, inner sep=1.5pt, font=\tiny\bfseries] (c0) at (s2-in-1.south) [yshift=-6pt] {0};
    \node[circle, draw=orange!90!black, fill=orange!40, inner sep=1.5pt, font=\tiny\bfseries] (c2) at (s2-in-3.south) [yshift=-6pt] {2};
    \node[circle, draw=orange!90!black, fill=orange!40, inner sep=1.5pt, font=\tiny\bfseries] (c4) at (s2-in-5.south) [yshift=-6pt] {4};

    % Arrow down to output (Starts cleanly below circle node c2)
    \draw[->, >=stealth, line width=1pt, black!60] (3.6, 1.7) -- (3.6, 1.1)
        node[midway, right, font=\scriptsize] {$K=3, P=0$};

    % Output Strip 1x3
    \node[out cell, fill=orange!40, draw=orange!80!black, line width=1.2pt] (s2-out-0) at (2.4, 0.4) {$a_0$};
    \node[out cell] (s2-out-1) at (3.6, 0.4) {$a_2$};
    \node[out cell] (s2-out-2) at (4.8, 0.4) {$a_4$};
    \node[font=\footnotesize\bfseries, color=green!50!black, anchor=east] at (-0.8, 0.4) {Output ($N_{\text{out}}=3$):};

    % Bottom Formula
    \node[labeltext, anchor=north] at (3.6, -0.55) {
        $N_{\text{out}} = \left\lfloor \frac{N_{\text{in}} + 2P - K}{S} \right\rfloor + 1 = \left\lfloor \frac{7 - 3}{2} \right\rfloor + 1 = \mathbf{3}$ \textcolor{green!50!black}{(\textbf{Downsampled})}
    };
\end{scope}

\end{tikzpicture}
\caption{Visualizing the effect of stride on spatial resolution ($N_{\text{in}}=7, K=3, P=0$). \textbf{Top:} Stride $S=1$ evaluates every valid position ($0,1,2,3,4$), producing 5 activations. \textbf{Bottom:} Stride $S=2$ skips intermediate starting locations ($1$ and $3$), evaluating only positions $0, 2, 4$ and downsampling the output to 3 activations.}
\label{fig:stride_comparison_vertical}
\end{figure}